We evaluated L6-guided Grover Adaptive Search for evidence selection after dense and lexical retrieval. A cross-encoder scores relevance, while a compact classical simulator applies phase marking and amplitude amplification to document subsets. On SciFact and NFCorpus, refinement kept the first five results and reduced mean redundancy relative to matched base selection. Recall changes and their paired uncertainty remain explicit: score constraints do not guarantee judged relevance. FiQA's refinement follow-up failed the recall gate, so its headline uses the earlier configuration. Greedy and exact controls matched the protected method's quality, supporting the selection policy without establishing a unique Grover benefit. Native circuit measurements show resource growth, with limits on preparation, simulation, and memory. Classical retrieval, neural scoring, and simulation remain part of the full cost. These findings support a reproducible relevance--redundancy tradeoff without establishing quantum speed or monetary savings. Future work should assess device connectivity, noise, energy use, and generated-answer quality.
